% plus3.tex -- Is AL+3 possible?  Landscape, structure of extensions of a projective plane,
% and computer evidence beyond the hat model.  Drop-in section: needs amsmath, amssymb, amsthm
% (theorem, lemma, proposition, remark environments), enumitem, booktabs.
% Code: research/plus3/ (sat_ext.py, jobs.py, thresholds.py, analyze.py); logs in runs/.
\providecommand{\PG}{\mathrm{PG}}
\providecommand{\cP}{\mathcal{P}}
\providecommand{\cL}{\mathcal{L}}
\providecommand{\cX}{\mathcal{X}}

\section{Beyond the hat model}\label{sec:plus3}

Write $n_b(\Delta,3)$ for the largest order of a bipartite graph of maximum degree $\Delta$ and diameter $3$.
For a prime power $q$, the construction of~\cite{AL} gives $n_b(q+2,3) \ge 2(q^2+q+5)$ and the four-hat
construction gives $n_b(q+2,3) \ge 2(q^2+q+6)$, two vertices more. Here we ask whether one more vertex is
possible, that is, whether $n_b(q+2,3) \ge 2(q^2+q+6)+1$ (``$+3$'' over~\cite{AL}).
We cannot settle this. We record where the question is open, what any such graph must look like when it is
built on a projective plane of order~$q$ and $q$ is large, and what a computer search covers.

\subsection{Where the question is open}\label{ssec:frontier}
The bipartite Moore bound for diameter $3$ is $2(\Delta^2-\Delta+1)$; for $\Delta = q+2$ it is $2(q^2+3q+3)$, which
exceeds $2(q^2+q+6)$ by $4q-6$. It is attained if and only if a projective plane of order $\Delta-1 = q+1$ exists.
When both $q$ and $q+1$ are prime powers (one of them is then a power of~$2$: $q \in \{2,3,4,7,8,16,31,\dots\}$),
$\PG(2,q+1)$ attains it and our question is void. Bipartite graphs of defect $1$ or $3$ do not exist
(for defect less than $\Delta$ such graphs are regular, hence of even order; see~\cite[Prop.~3.1]{FPPV}),
and the defect-$2$ case at diameter $3$ is restricted by~\cite{DJMP}. None of these upper bounds comes near:
for $q \ge 9$ the value $2(q^2+q+6)+1$ has defect $4q-7 \ge 29$.

In the range of the Combinatorics Wiki table ($\Delta \le 16$) the relevant cells are:
\begin{center}
\begin{tabular}{@{}rrrrrl@{}}
\toprule
$\Delta$ & $q$ & \cite{AL} & four hats & table upper bound & status of $2(q^2+q+6)+1$\\
\midrule
$7$ & $5$ & $70$ & $72$ & $80$ & below $n_b(7,3) = 80$ (Hafner--Loz graph; no $(7,3,-4)$-graph~\cite{FMPV})\\
$11$ & $9$ & $190$ & $192$ & $220$ & open ($193$)\\
$13$ & $11$ & $274$ & $276$ & $314$ & open ($277$)\\
$15$ & $13$ & $374$ & $376$ & $422$ & open ($377$)\\
\bottomrule
\end{tabular}
\end{center}
For $\Delta = 16$ ($q+1 = 15$) the table value $394$ comes from duplication and is not of this form.
Beyond the table, $2(q^2+q+6)+1$ is the open value for every prime power $q \ge 9$ with $q+1$ not a prime power
($q = 9, 11, 13, 17, 19, 23, 25, 27, 29, 32, 37, \dots$), as far as we know.

\subsection{Extensions of a projective plane}\label{ssec:ext}
Every graph in \S\ref{sec:coords} contains the incidence graph of $\PG(2,q)$ with a few incidences changed.
We now ask what any graph of this kind must look like.

Let $\Pi$ be a projective plane of order $q$ with point set $\cP$ and line set $\cL$, and $N = q^2+q+1$.
An \emph{extension} of $\Pi$ is a bipartite graph $G$ with parts $\cP \cup P'$ (points) and $\cL \cup L'$ (lines),
of maximum degree at most $q+2$ and diameter at most $3$. The vertices in $P'$ and $L'$ are \emph{new}, the others
\emph{old}. A \emph{change} is a pair (old point, old line) on which $G$ and $\Pi$ disagree: a \emph{deleted} flag
or an \emph{added} non-incident pair. Write $p = |P'|$, $l = |L'|$ and $e$ for the number of changes.
The order of $G$ is $2N+p+l$; so~\cite{AL} has $p = l = 4$, the four-hat graph has $p = l = 5$, and ``$+3$'' asks for $p+l = 11$.
In a bipartite graph without isolated vertices, diameter at most $3$ means: any two points have a common neighbour
and any two lines have a common neighbour (a point $x$ and a line $L \not\ni x$ are then at distance $3$ through
any point of $L$).

For a new point $h$, let $O(h)$ be the set of old lines through $h$, $a_h = |O(h)|$, and $d_h$ the number of new lines
through $h$. Dually $O(K)$, $c_K$ for a new line $K$. For an old point $x$ let $r_x(h)$ be the number of lines of
$O(h)$ that pass through $x$ \emph{in $\Pi$}, and $\sigma_h = \sum_x \max(r_x(h)-1, 0)$.

\begin{lemma}\label{lem:budget}
\begin{enumerate}[label=\textup{(\alph*)}]
\item An old point $x$ lies on at most $1 + \delta(x) - \alpha(x)$ new lines, where $\delta(x)$ and $\alpha(x)$ count
the deleted and added pairs at $x$. Dually for old lines.
\item Two new points with no common new line lie on a common old line carrying a deleted flag.
In particular, if $e = 0$, any two new points lie on a common new line, and any two new lines meet in a new point.
\end{enumerate}
\end{lemma}
\begin{proof}
(a) $x$ has $q+1-\delta(x)+\alpha(x)$ old neighbours and degree at most $q+2$.
(b) By (a), an old line through two new points has a deleted flag.
\end{proof}

\begin{lemma}\label{lem:cover}
Let $h$ be a new point, $a = a_h$, $d = d_h$, and let $t$ be the number of added pairs $(x,L)$ with $L \in O(h)$.
Then $\sigma_h \le (a+d)(q+1) - N + t \le 2q+1+t$, and $a+d \ge q+1-\varepsilon$, where $\varepsilon = 0$ if $t = 0$ and
$\varepsilon = 1$ if $1 \le t \le q+1$.
\end{lemma}
\begin{proof}
In $\Pi$ the $a$ lines of $O(h)$ cover $a(q+1)-\sigma_h$ points; in $G$ they cover at most $t$ more.
Every other old point is at distance $2$ from $h$ through a new line, and a new line through $h$ has at most $q+1$ old
points. Hence $N - a(q+1) + \sigma_h - t \le d(q+1)$. Use $a + d \le q+2$ for the first claim and $\sigma_h \ge 0$ for the second.
\end{proof}

\begin{lemma}[near-copies]\label{lem:copy}
Let $h$, $a$, $t$, $\varepsilon$ be as in Lemma~\ref{lem:cover} and suppose $q \ge 3l + 6 + t + 3\varepsilon$. Then there is an
old point $X$ such that all but at most one line of $O(h)$ pass through $X$ (in $\Pi$). At least $q-l-\varepsilon-1$ lines of
$O(h)$ pass through $X$.
\end{lemma}
\begin{proof}
Let $X$ maximise $r = r_X(h)$, and let $b = a - r$ be the number of lines of $O(h)$ not through $X$; put $s = 2q+1+t$,
so $\sigma_h \le s$ by Lemma~\ref{lem:cover}. Two facts:
(i) $(b+1)r \le s+1$. Indeed $X$ contributes $r-1$ to $\sigma_h$; every point $x \ne X$ on one of the $r$ lines through $X$
lies on exactly one of them and contributes the number of non-$X$ lines through it; and each of the $b$ non-$X$
lines meets each of the $r$ lines in a point other than $X$. So $\sigma_h \ge r - 1 + br$.
(ii) $a(a-1) \le rs$. Any two lines of $O(h)$ meet in one point of $\Pi$, so
$a(a-1) = \sum_x r_x(r_x-1) \le r\sum_x \max(r_x-1,0) = r\sigma_h$.

By Lemma~\ref{lem:cover}, $a \ge q+1-l-\varepsilon$, and the hypothesis gives $a \ge (2q+9+t)/3 = (s+8)/3$. Suppose $b \ge 2$.
Then $r \le (s+1)/3$ by (i), $r \ge r_0 := a(a-1)/s$ by (ii), and $a = r + b \le g(r) := r + (s+1)/r - 1$ by (i).
As $g$ is convex, $a \le \max\{g(r_0), g((s+1)/3)\}$. First, $g((s+1)/3) = (s+7)/3 < a$. Second, $a/3 < r_0 \le a(q+1)/(2q+1)$
(use $a \ge (s+5)/3 + 1$ and $a \le q+2$), and $x \mapsto x + (s+1)/x$ is convex, so $g(r_0) + 1$ is at most the larger of
$a/3 + 3(s+1)/a < a/3 + 9$ and $a(q+1)/(2q+1) + (s+1)(2q+1)/(a(q+1)) < a(q+1)/(2q+1) + 6$. Both are at most $a+1$ once
$a \ge 13$, which holds for $q \ge 15$. So $g(r_0) < a$, a contradiction. For $q \le 14$ the hypothesis forces $l \le 2$ and
$t \le 5$, and an exhaustive check of the integer system (i), (ii), $a = r + b$, $b \ge 2$, $a \ge q+1-l-\varepsilon$ finds no
solution (\texttt{thresholds.py}, which checks all $q \le 400$, $0 \le l \le 12$, $t \le 12$; the bound $3l+6+t+3\varepsilon$ is
exactly the threshold of this system in every case checked). Hence $b \le 1$ and $r = a - b \ge q - l - \varepsilon - 1$.
\end{proof}

We call $X$ the point \emph{copied} by $h$ and write $X = X_h$; dually a new line $K$ copies an old line $L_K$.
In the four-hat graph, $Q$ copies $A$, $H_s$ copies $X_s$, $M$ copies $\ell$ and $K_s$ copies $AX_s$.

\begin{theorem}\label{thm:copies}
Let $G$ be an extension of $\Pi$ with $p$ new points, $l$ new lines and $e$ changes, and suppose
$q \ge 3\max(p,l) + 9 + e$. Then:
\begin{enumerate}[label=\textup{(\alph*)}]
\item every new point $h$ copies an old point $X_h$: all but at most one of its old lines pass through $X_h$, and it
has at least $q - l - 2$ old lines through $X_h$; dually every new line $K$ copies an old line $L_K$;
\item distinct new points copy distinct points, and distinct new lines copy distinct lines;
\item for every line $M$ of $\Pi$ through $X_h$ that is not an old line through $h$, some new line $K$ through $h$ copies $M$.
If $m_h$ is the number of such lines $M$, then $m_h \le d_h \le m_h + 1$, and $d_h = m_h$ if $h$ lies on an old line not
through $X_h$. Dually for new lines.
\end{enumerate}
\end{theorem}
\begin{proof}
(a) is Lemma~\ref{lem:copy} with $t \le e$ and $\varepsilon \le 1$, and its dual.
(b) If $X_h = X_{h'}$ for new points $h \ne h'$, then at least $2(q-l-2) - (q+1)$ lines through $X_h$ are old lines
through both $h$ and $h'$; each of them carries a deleted flag by Lemma~\ref{lem:budget}(a), so $q - 2l - 5 \le e$,
against the hypothesis. Dually for lines.
(c) The $r$ old lines of $h$ through $X_h$ meet $M$ only in $X_h$, the possible extra old line of $h$ meets it in one
point, and added pairs put at most $e$ further points of $M$ on old lines of $h$. So at least $q-1-e$ points of $M$ are
at distance $2$ from $h$ only through new lines through $h$. A new line $K$ through $h$ with $L_K \ne M$ has at most two
old points on $M$ (one on $L_K$, and possibly its one old point off $L_K$). If no $K$ through $h$ copied $M$, then
$q-1-e \le 2d_h \le 2l$, which is false. The copies of different lines $M$ are different by (b), so $d_h \ge m_h$.
Finally $a_h = (q+1-m_h) + b$ with $b \in \{0,1\}$ the number of old lines of $h$ not through $X_h$, and $a_h + d_h \le q+2$.
\end{proof}

So, for large $q$, every graph of order $2N+k$ built on a projective plane by $e$ changes is a \emph{copy
configuration}: a set $\cX$ of copied points and a set $\cL'$ of copied lines, such that every line through a copied point
is either an old line of its copy or is itself copied (and that copy passes through the point's copy), and dually.
Each new vertex has at most one ``free'' new neighbour that is not explained by this rule. This is the
general form of the hat lemma (\cite{AL}; \S\ref{sec:coords}): in the four-hat graph $\cX = \{A, X_0, X_1, X_a, X_c\}$
and $\cL' = \{\ell, AX_0, AX_1, AX_a, AX_c\}$; the hats $H_0, H_1, H_c$ have $m = 1$ (their missing line $\ell$ is
copied by $M$), $H_a$ lies on $\ell$ and has $m = 0$, and $Q$ has $m = 3$ (missing $AX_0$, $AX_1$, $AX_c$, copied by
$K_0$, $K_1$, $K_c$; its free line is $K_a$). The hat model of the SAT search in \S\ref{sec:coords} fixes this
particular configuration. Theorem~\ref{thm:copies} says that for large $q$ only copy configurations occur; the next
proposition shows that without changes the configuration must have the shape of the hat configuration.

\begin{proposition}\label{prop:closed}
Let $G$ be an extension of $\Pi$ with no changes ($e = 0$), $q \ge 3\max(p,l)+9$ and $p + l \le 13$. Then the line
joining any two copied points is copied, and the point where any two copied lines meet is copied. Consequently all
copied points but at most one lie on one line, and all copied lines but at most one pass through one point.
\end{proposition}
\begin{proof}
Let $h \ne h'$ be new points and $J = X_hX_{h'}$ (the points are distinct by Theorem~\ref{thm:copies}(b)). An old line
has at most one new point when $e = 0$ (Lemma~\ref{lem:budget}(a)), so $J$ is not an old line of both $h$ and $h'$;
if it is not an old line of $h$, Theorem~\ref{thm:copies}(c) says $J$ is copied. The second claim is dual.
So $(\cX, \cL')$ is closed under joins and meets. If $\cX$ contained four points with no three collinear, the closure
would be a subplane, with at least $7$ points and $7$ lines, against $p + l \le 13$. A point set without such a
quadrangle has all but at most one of its points on one line: otherwise take a line $L$ with the most points of the
set, $m \ge 2$ of them, and two points $u, v$ of the set off $L$; if $m \ge 3$, two points of $L$ off the line $uv$
together with $u, v$ form a quadrangle, and if $m = 2$ any four points do. Dually for $\cL'$.
\end{proof}

With changes the same argument shows that at most $e$ joins of copied points can fail to be copied (each failure needs
a deleted flag on that line), so the configuration is closed up to $e$ exceptions. In the hat configuration the line
is $\ell$ and the point is $A$; the copied point off $\ell$ is $A$ and the copied line not through $A$ is $\ell$.

\begin{theorem}\label{thm:nochange}
Let $G$ be an extension of $\Pi$ with no changes, $p$ new points and $l$ new lines, $p + l \le 13$ and
$q \ge 3\max(p,l)+9$. Then $p \le 6$ and $l \le 6$, and $p = 6$ or $l = 6$ forces $p = l = 6$.
In particular $p + l \notin \{11, 13\}$: for $q \ge 3(p+l)+9$, hence for $q \ge 42$, no graph of maximum degree $q+2$
and diameter $3$ consists of the incidence graph of a projective plane of order $q$ as an induced subgraph plus exactly
$11$ (or $13$) further vertices.
\end{theorem}
\begin{proof}
By Proposition~\ref{prop:closed} there is a line $L^*$ containing all copied points but at most one; let
$S = \cX \cap L^*$ and $s = |S| \ge p-1$. Suppose $s \ge 6$. Then $L^*$ is copied (it joins two copied points),
say by $K^*$; write $h_Y$ for the new point copying $Y \in S$.
\begin{enumerate}[label=(\roman*)]
\item $L^*$ has at most one new point (Lemma~\ref{lem:budget}), so it is an old line of $h_Y$ for at most one $Y \in S$.
For every other $Y$, $L^*$ is missing at $h_Y$, so $K^*$ passes through $h_Y$ (Theorem~\ref{thm:copies}(c)).
\item The new points on $K^*$ are the copies of the points of $L^*$ that are not old points of $K^*$, and at most one
more (the dual of Theorem~\ref{thm:copies}(c)). Hence at most two points $Y \in S$ are old points of $K^*$: the one with
$L^* \in O(h_Y)$, and one whose copy is the extra new point of $K^*$.
\item A new line $K \ne K^*$ contains at most two copies $h_Y$, $Y \in S$: its line $L_K$ meets $L^*$ in one point, and
$K$ has at most one new point that is not the copy of a point of $L_K$.
\item Let $Y \ne Y'$ in $S$ with $L^*$ missing at $h_{Y'}$. The old lines of $h_{Y'}$ are lines through $Y'$ other than
$L^* = YY'$, and at most one further line $E_{Y'}$. So unless $Y \in E_{Y'}$, the points $Y$ and $h_{Y'}$ are joined only
by a new line, which must be the unique new line $\nu(Y)$ through $Y$.
\end{enumerate}
Let $e_Y$ be the number of $Y' \in S$ with $Y \in E_{Y'}$; as each $E_{Y'}$ meets $L^*$ once, $\sum_Y e_Y \le s$.
By (i) and (iv), at least $s - 2 - e_Y$ copies $h_{Y'}$ must lie on the new line $\nu(Y)$ (if this number is positive,
$\nu(Y)$ exists). For the at least $s-2$ points $Y$ that are not old
points of $K^*$ (by (ii)), $\nu(Y) \ne K^*$, so by (iii) $s - 2 - e_Y \le 2$, that is $e_Y \ge s-4$ (also when $\nu(Y)$ does not exist). Then
$(s-2)(s-4) \le \sum_Y e_Y \le s$, which fails for $s \ge 6$. So $s \le 5$ and $p \le 6$.
If $p = 6$, then $s = 5$ and there is a copied point $u \notin L^*$; the five lines $uY$, $Y \in S$, are copied
(Proposition~\ref{prop:closed}) and so is $L^*$, so $l \ge 6$. The statements for $l$ are dual. If $p + l \in \{11,13\}$
then $\max(p,l) \ge 6$ and $\min(p,l) \le 6$ with $(p,l) \ne (6,6)$, a contradiction.
\end{proof}

The hypothesis $q \ge 3\max(p,l)+9$ is what Theorem~\ref{thm:copies} needs; the counting in the proof itself does not
depend on $q$. With changes the argument breaks in two places: a changed incidence can let two new points share an
old line, and can let a copied point see a hat through an added old incidence. So a $+3$ on a projective plane of large
order must change the plane, and Theorem~\ref{thm:nochange} says nothing about how many changes would suffice.
Whether $p = l = 6$ (that is, $2(q^2+q+7)$, five hats) can occur without changes is also left open by the counting;
the hat-model SAT search of \S\ref{sec:coords} found no such graph.

\subsection{The four hats without changing the plane}\label{ssec:zero}
The four-hat construction of \S\ref{sec:coords} moves three incidences of the plane (the rotation). The computer
search below found that this is not necessary: the rotation can be carried by the copies.

\begin{proposition}\label{prop:zero}
Let $\Pi$ be a projective plane of order $q \ge 3$, $\ell$ a line, $A$ a point not on $\ell$, and $X_a, X_1, X_2, X_3$
distinct points of $\ell$. Let $s$ be the $3$-cycle $1 \mapsto 2 \mapsto 3 \mapsto 1$ and write $AX$ for the line through $A$
and $X$. Add new points $Q, H_a, H_1, H_2, H_3$ and new lines $M, K_a, K_1, K_2, K_3$ to the incidence graph of $\Pi$,
with these incidences ($i \in \{1,2,3\}$):
\begin{center}
\begin{tabular}{@{}l@{\quad}l@{\quad}l@{}}
\toprule
vertex & old neighbours & new neighbours\\
\midrule
$Q$ & every line through $A$ except $AX_1, AX_2, AX_3$ & $K_a, K_1, K_2, K_3$\\
$H_a$ & every line through $X_a$ except $AX_a$ & $M, K_a$\\
$H_i$ & every line through $X_i$ except $\ell$ and $AX_i$; also $AX_{s(i)}$ & $M, K_i$\\
$M$ & every point of $\ell$ except $X_1, X_2, X_3$ & $H_a, H_1, H_2, H_3$\\
$K_a$ & every point of $AX_a$ except $X_a$ & $Q, H_a$\\
$K_i$ & every point of $AX_i$ except $A$ and $X_i$; also $X_{s^{-1}(i)}$ & $Q, H_i$\\
\bottomrule
\end{tabular}
\end{center}
The resulting graph $G$ is bipartite, has maximum degree $q+2$, diameter $3$ and order $2(q^2+q+6)$, and every
incidence of $\Pi$ is kept: the incidence graph of $\Pi$ is an induced subgraph of $G$.
\end{proposition}

\begin{proof}
\emph{Degrees.} New vertices: $Q$ has $(q+1-3)+4$, $H_a$ has $q+2$, $H_i$ has $(q-1)+1+2$, $M$ has $(q-2)+4$, $K_a$ has
$q+2$ and $K_i$ has $(q-1)+1+2$ neighbours. Every old vertex gets at most one new neighbour: a line through $A$ meets
$\ell$ in one point, so the old lines with a new point are $AX_a$ and the lines $AP$, $P \in \ell\setminus\{X_a,X_1,X_2,X_3\}$
(point $Q$), the lines through $X_a$ other than $AX_a$, among them $\ell$ (point $H_a$), the lines through $X_i$ other than
$\ell$ and $AX_i$ (point $H_i$), and $AX_i$ (point $H_{s^{-1}(i)}$), and these classes are disjoint. Dually the old
points with a new line are the points of $\ell \setminus \{X_1,X_2,X_3\}$ (line $M$), of $AX_a \setminus \{X_a\}$ (line $K_a$),
of $AX_i \setminus\{A, X_i\}$ (line $K_i$), and $X_i$ (line $K_{s(i)}$).

\emph{Two points.} Old pairs are joined in $\Pi$. $Q$ and an old point $x$: if $Ax \ne AX_i$ for all $i$ then $Ax$ is a line
of $Q$; if $x \in AX_i \setminus \{A, X_i\}$, use $K_i$; if $x = X_i$, use $K_{s(i)}$. $H_a$ and $x$: if $x \notin AX_a$ or
$x = X_a$ use $X_ax$; otherwise use $K_a$. $H_i$ and $x$: if $X_ix \notin \{\ell, AX_i\}$ use it; $A$ and $X_{s(i)}$ lie on
$AX_{s(i)}$; the points of $\ell$ other than $X_1,X_2,X_3$ lie on $M$; $X_{s^{-1}(i)}$ and the points of
$AX_i \setminus\{A,X_i\}$ lie on $K_i$; and $X_i$ is on every line of $H_i$ through it. New pairs: $Q, H_a$ on $K_a$;
$Q, H_i$ on $K_i$; any two hats on $M$.

\emph{Two lines.} $M$ and an old line $L$: if $L$ meets $\ell$ in a point of $M$, done; if $L \ni X_i$ and $L \ne \ell$, then
$L$ contains $H_i$ (when $L \ne AX_i$) or $H_{s^{-1}(i)}$ (when $L = AX_i$), and both lie on $M$. $K_a$ and $L$: $L$ meets
$AX_a$; if only in $X_a$, then $L$ is a line through $X_a$ other than $AX_a$ and contains $H_a$. $K_i$ and $L \ne AX_i$:
let $y = L \cap AX_i$. If $y \notin \{A, X_i\}$, done. If $y = X_i$ and $L \ne \ell$, $L$ contains $H_i$; if $L = \ell$, it
contains $X_{s^{-1}(i)} \in K_i$. If $y = A$, then $L$ contains $Q$ unless $L = AX_j$ with $j \ne i$; then either
$j = s(i)$ and $L$ contains $H_i$, or $j = s^{-1}(i)$ and $L$ contains $X_j \in K_i$. New pairs: $M$ meets $K_a$ in $H_a$
and $K_i$ in $H_i$; the $K$'s meet in $Q$.
\end{proof}

The program \texttt{construction0.py} builds $G$ from the table alone and checks it on random frames for
$q = 3,4,5,7,8,9,11,13,16$ and on the nearfield plane of order $9$; the graphs for $q = 4, 5, 9, 11, 13$ pass
\texttt{verify\_record.py}, and removing either rotation incidence ($H_i \sim AX_{s(i)}$ or $X_{s^{-1}(i)} \sim K_i$)
breaks the diameter. In the language of Theorem~\ref{thm:copies}: $H_1, H_2, H_3$ each miss two lines through their point
($\ell$ and $AX_i$, copied by $M$ and $K_i$) and carry one old line not through it ($AX_{s(i)}$); $K_1, K_2, K_3$ are dual.
An AL-size graph without changes ($p = l = 4$: $Q$, $M$ and three hats whose rotation is a $3$-cycle; $Q$ and $M$
each carry one extra old incidence) was found by the same search at $q = 7$ (\texttt{graphs/ext\_q7\_4\_4\_E0\_A0\_n122.json}).
